\documentclass[11pt,a4paper]{article}
\usepackage{fontspec,amsmath,amssymb,amsthm,microtype}
\setmainfont{lmroman10-regular.otf}[BoldFont=lmroman10-bold.otf,ItalicFont=lmroman10-italic.otf,BoldItalicFont=lmroman10-bolditalic.otf]
\usepackage{xcolor}
\usepackage[most]{tcolorbox}
\usepackage{mathrsfs}
\providecommand{\lt}{<}
\providecommand{\gt}{>}
\setmonofont{DejaVuSansMono.ttf}[BoldFont=DejaVuSansMono-Bold.ttf,ItalicFont=DejaVuSansMono-Oblique.ttf,BoldItalicFont=DejaVuSansMono-BoldOblique.ttf,Scale=0.86]
\newfontfamily\leanfallback{STIXTwoMath-Regular.otf}[Scale=0.86]
\definecolor{leanrule}{gray}{0.62}
\definecolor{leanback}{gray}{0.962}
\definecolor{leantitle}{gray}{0.22}
\definecolor{leanmuted}{gray}{0.38}
\newtcolorbox{leanbox}[1]{enhanced,breakable,colback=leanback,colframe=leanrule,
  boxrule=0.4pt,arc=0.6mm,left=2.4mm,right=2.4mm,top=1.8mm,bottom=1.8mm,
  before skip=0.8em,after skip=1em,fonttitle=\small\bfseries,coltitle=leantitle,
  colbacktitle=leanback,titlerule=0.3pt,title={#1}}
\newlength{\leancharwidth}
\newenvironment{leancode}{\par\smallskip\begingroup\ttfamily\footnotesize\raggedright
  \settowidth{\leancharwidth}{M}\setlength{\parskip}{0pt}\setlength{\parindent}{0pt}
  \frenchspacing}{\par\endgroup\smallskip}
\newcommand{\leanline}[2]{\hangindent=\dimexpr#1\leancharwidth+4\leancharwidth\relax
  \hangafter=1\noindent\hspace*{#1\leancharwidth}#2\par}
\newcommand{\leansource}[1]{\par{\footnotesize\color{leanmuted}#1}\par}
\newcommand{\leanlabel}[1]{\par\smallskip\noindent{\small\bfseries #1}\par}
\usepackage[margin=25mm]{geometry}
\usepackage[unicode,hidelinks]{hyperref}
\hypersetup{pdfauthor={LangevinLean Project}}
\setlength{\emergencystretch}{3em}
\setlength{\parskip}{0.4em}
\setlength{\parindent}{0pt}
\allowdisplaybreaks[2]
\raggedbottom
\title{\texorpdfstring{Self-adjointness of the actual Langevin transition}{Self-adjointness of the actual Langevin transition}}
\author{LangevinLean Project}
\date{}
\begin{document}\maketitle
Draft mathematical exposition; these notes may contain mistakes.


Actual anchored random MALA pairs are reversible. Their corresponding Gibbs observable compressions converge strongly to the actual Langevin operator. Passing the symmetric pairing identity to that genuine limit proves actual diffusion reversibility.


\section{\texorpdfstring{The true reversible approximants and their diffusion limit}{The true reversible approximants and their diffusion limit}}


\textbf{Claim.} The actual canonical Langevin operator \(P_t\) is self-adjoint for \(t\ge0\).



\textbf{Proof.} The zero-time diffusion variable is the actual initial coordinate. The exact anchored sampled MALA pair law therefore identifies its initial/\allowbreak{}endpoint pair with the reversible explicit kernel-power pair. The pair-law theorem makes each actual MALA compression self-adjoint. Choose the floor-mesh counts approaching \(t\). The actual MALA operator approximation theorem gives strong convergence to \(P_t\); self-adjointness passes to that strong limit by testing all pairings against fixed vectors.

\begin{leanbox}{Lean: \texttt{diffusion\_\allowbreak{}zero\_\allowbreak{}coe}, \texttt{mala\_\allowbreak{}pair\_\allowbreak{}swap}, \texttt{malaOp\_\allowbreak{}isSelfAdjoint} and 1 more}
\leansource{\texttt{diffusion\_\allowbreak{}zero\_\allowbreak{}coe} \quad WassersteinMixing/\allowbreak{}SemigroupReversible.lean, lines 21–28}
\begin{leancode}
\leanline{0}{theorem\ diffusion\_zero\_coe\ (P\ :\ Parameters)\ (V\ :\ E\ P.dimension\ →\ ℝ)}
\leanline{4}{(hV\ :\ PotentialAssumptions\ P\ V)\ :}
\leanline{4}{diffusion\ P\ V\ hV\ 0\ =ᵐ[probability\ P\ V]\ fun\ ω\ ↦\ actualInitial\ P\ V\ hV\ ω.1}
\end{leancode}
\leansource{\texttt{mala\_\allowbreak{}pair\_\allowbreak{}swap} \quad WassersteinMixing/\allowbreak{}SemigroupReversible.lean, lines 32–46}
\begin{leancode}
\leanline{0}{theorem\ mala\_pair\_swap\ (P\ :\ Parameters)\ (V\ :\ E\ P.dimension\ →\ ℝ)}
\leanline{4}{(hV\ :\ PotentialAssumptions\ P\ V)\ (r\ :\ ℝ)\ (hr\ :\ 0\ <\ r)\ (n\ :\ ℕ)\ :}
\leanline{4}{((probability\ P\ V).map\ (fun\ ω\ ↦}
\leanline{6}{(diffusion\ P\ V\ hV\ 0\ ω,\ mala\ P\ V\ hV\ r\ hr\ n\ ω))).map\ Prod.swap\ =}
\leanline{6}{(probability\ P\ V).map\ (fun\ ω\ ↦}
\leanline{8}{(diffusion\ P\ V\ hV\ 0\ ω,\ mala\ P\ V\ hV\ r\ hr\ n\ ω))}
\end{leancode}
\leansource{\texttt{malaOp\_\allowbreak{}isSelfAdjoint} \quad WassersteinMixing/\allowbreak{}SemigroupReversible.lean, lines 53–59}
\begin{leancode}
\leanline{0}{theorem\ malaOp\_isSelfAdjoint\ (P\ :\ Parameters)\ (V\ :\ E\ P.dimension\ →\ ℝ)}
\leanline{4}{(hV\ :\ PotentialAssumptions\ P\ V)\ (r\ :\ ℝ)\ (hr\ :\ 0\ <\ r)\ (n\ :\ ℕ)\ :}
\leanline{4}{IsSelfAdjoint\ (malaOp\ ({\leanfallback 𝕜}\ :=\ {\leanfallback 𝕜})\ (F\ :=\ F)\ P\ V\ hV\ r\ hr\ n)}
\end{leancode}
\leansource{\texttt{op\_\allowbreak{}isSelfAdjoint} \quad WassersteinMixing/\allowbreak{}SemigroupReversible.lean, lines 70–87}
\begin{leancode}
\leanline{0}{theorem\ op\_isSelfAdjoint\ (P\ :\ Parameters)\ (V\ :\ E\ P.dimension\ →\ ℝ)}
\leanline{4}{(hV\ :\ PotentialAssumptions\ P\ V)\ (t\ :\ ℝ)\ (ht\ :\ 0\ ≤\ t)\ :}
\leanline{4}{IsSelfAdjoint\ (op\ ({\leanfallback 𝕜}\ :=\ {\leanfallback 𝕜})\ (F\ :=\ F)\ P\ V\ hV\ t)}
\end{leancode}
\leanlabel{How the Lean proof follows the argument}
\texttt{mala\_\allowbreak{}pair\_\allowbreak{}swap} uses the exact canonical anchored law and actual initial representative. \texttt{mala\allowbreak{}Op\_\allowbreak{}is\allowbreak{}Self\allowbreak{}Adjoint} applies pair-law reversal. \texttt{op\_\allowbreak{}is\allowbreak{}Self\allowbreak{}Adjoint} combines the floor-mesh convergence, actual operator approximation and the proved strong-limit self-adjointness theorem.
\end{leanbox}


Earlier results used here: \cite[Section~3.26.3]{LangevinLean2026ULASC}.

\begin{thebibliography}{ULA-LSI}
\bibitem[ULA-SC]{LangevinLean2026ULASC}
LangevinLean Project. \emph{Unadjusted Langevin under Strong Convexity}. LL-01, 2026. Draft report, content revision d34ead06987c. \url{https://odakyildiz.github.io/LangevinLean/docs/pdf/UnadjustedLangevin/FullProof.pdf}.
\end{thebibliography}

\end{document}
